\begin{table*}[!t]
\hfuzz=2pt
\centering
\begin{minipage}{\linewidth}
\section{Illustrative Examples}
\label{app:rationale-analysis}
Tables~\ref{tab:rationale-quality-example}--\ref{tab:rationale-correction-example}
retain the task, gold label, score-determining input and rationale evidence,
predictions, and audit outcomes while omitting repeated full outputs.
Tables~\ref{tab:rationale-quality-example} and~\ref{tab:rationale-consistency-example} contrast external quality preference with support for each model's own score. Tables~\ref{tab:rationale-sentence-example} and~\ref{tab:rationale-correction-example} trace a rubric error through correction and fixed-prefix rescoring.
\end{minipage}

\vspace{0.7\baselineskip}
\begin{minipage}{\textwidth}
\hfuzz=2pt
\normalsize\linespread{1.05}\selectfont
\renewcommand{\arraystretch}{1.08}
\textbf{Task: Grounding Specificity}\hfill\textbf{Gold: 3}\\[2pt]
\begin{tabularx}{0.95\linewidth}{@{}>{\raggedright\arraybackslash}p{0.22\linewidth}X@{}}
\toprule
Evaluation input & \textit{``Limitations of modeling with GP: We know that the
required number of samples to accurately model the objective function grows
exponentially, leading to increased computational complexity and potentially
reduced accuracy. Do you plan to address this limitation in future work?''} \\
\midrule
SO rationale (pred. 1) & Treats the comment as a general GP limitation and
concludes that it does not identify a specific part or paper-specific issue. \\
RS rationale (pred. 3) & Distinguishes grounding from specificity: the paper
location is imprecise, but exponential sample growth and its consequences form
a clearly specified issue. \\
\bottomrule
\end{tabularx}
\vspace{2pt}
{\normalsize
\setlength{\tabcolsep}{0.6pt}
\begin{tabular*}{0.95\linewidth}{@{\extracolsep{\fill}}lrrrrl@{}}
& Rubric & Evidence & Coverage & Unsup. & Pref. \\
SO & 1.0 & 1.0 & 1.5 & 1.0 & \\
RS & 3.0 & 3.0 & 3.0 & 3.0 & \textbf{RS} \\
\end{tabular*}}
\captionsetup{hypcap=false,width=0.96\linewidth}
\captionof{table}{Rationale-quality example. The four columns report rubric
correctness, evidence grounding, decision coverage, and unsupported-claim
control; the aggregated preference is RS.}
\label{tab:rationale-quality-example}
\end{minipage}
\vspace{0.8\baselineskip}
\begin{minipage}{\textwidth}
\normalsize\linespread{1.05}\selectfont
\setlength{\tabcolsep}{7pt}
\renewcommand{\arraystretch}{1.12}
\begin{tabularx}{\linewidth}{@{}X|X@{}}
\toprule
\multicolumn{2}{@{}p{\dimexpr\linewidth-2\tabcolsep-\arrayrulewidth\relax}@{}}{%
\textbf{Task: Helpfulness} \quad $\vert$ \quad
\textbf{Evaluation: Rationale--Score Consistency Assessment} \hfill
\textbf{Gold label $=5$}} \\
\midrule
\multicolumn{2}{@{}>{\raggedright\arraybackslash}p{\dimexpr\linewidth-2\tabcolsep-\arrayrulewidth\relax}@{}}{%
\textbf{Evaluation input:} \textit{Overall training time. The overall training
time and the bandwidth of exchanging hidden vectors should be mentioned. If the
network size of the exchanged information is too high and bandwidth is limited
in wireless communication, a single iteration will take too much time or fail
because of communication-protocol constraints. For example, in an IoT setting,
the maximum payload length is only 65,535 bytes. Please discuss the size of the
exchanged information (hidden vectors or gradients) and the total training time
in the revision.}} \\
\midrule
\rowcolor{gray!12}
\textbf{SO\textsubscript{RF} rationale} \hfill \textbf{predicts 5} &
\textbf{RS\textsubscript{RF} rationale} \hfill \textbf{predicts 4} \\
Calls the feedback comprehensive and detailed because it identifies
training-time and bandwidth constraints and requests concrete reporting. &
Recognizes that the request is clear and actionable, but could provide more
detail about how to measure or report the quantities. \\
\end{tabularx}
\nopagebreak
\vspace{-2pt}
\begin{tabularx}{\linewidth}{@{}l
    >{\centering\arraybackslash}p{0.22\linewidth}
    >{\raggedright\arraybackslash}X@{}}
\midrule
 & \textbf{Verdict} & \textbf{Judge explanation} \\
\textbf{Judge on SO} & \texttt{partially\_}\allowbreak\texttt{supported} &
The rationale identifies genuine strengths but overstates breadth and
implementation guidance. \\
\rowcolor{gray!12}
\textbf{Judge on RS} & \textbf{\texttt{supported}} &
The stated limitation supports the model's own score; overall preference is RS. \\
\midrule
\multicolumn{3}{@{}p{\dimexpr\linewidth-2\tabcolsep\relax}@{}}{%
\hfill \textbf{Overall preference: \underline{RS}.}} \\
\bottomrule
\end{tabularx}
\captionsetup{hypcap=false}
\captionof{table}{Internal rationale--score consistency example. The verdicts
assess support for each model's own prediction, not correctness against gold.}
\label{tab:rationale-consistency-example}
\end{minipage}
\end{table*}

\begin{table*}[!t]
\hfuzz=2pt
\centering
\begin{minipage}{\textwidth}
\normalsize\linespread{1.05}\selectfont
\renewcommand{\arraystretch}{1.08}
\textbf{Task: Actionability}\hfill
\textbf{SO seed 42; Gold/DS/RF: 3/3/5}\\[2pt]
\begin{tabularx}{0.99\linewidth}{@{}>{\raggedright\arraybackslash}p{0.22\linewidth}X@{}}
\toprule
Evaluation input & \textit{The comment questions the motivation for a
convolutional-coding-based DUN because its coefficients and dictionary are
learned end to end, argues that added hidden-layer channels may provide the same
capacity, and asks the authors to give more convincing motivation, e.g., a
theorem or experiments.} \\
\midrule
RF rationale & ``The answer provides a specific suggestion for the authors to
`give more convincing motivation (e.g., theorem, experiments).' This is an
explicit action with concrete details on how to implement it. The authors can
directly identify what needs to be done and how to apply the action. Therefore,
this meets the criteria for a score of 5.'' \\
\bottomrule
\end{tabularx}
\vspace{2pt}
\begin{tabularx}{\linewidth}{@{}>{\raggedright\arraybackslash}p{0.37\linewidth}>{\raggedright\arraybackslash}p{0.25\linewidth}X@{}}
\toprule
Flagged sentence & Error type & Judge basis \\
\midrule
``This is an explicit action with concrete details on how to implement it.'' &
Rubric misapplication & Naming evidence categories does not provide concrete
implementation details. \\
``The authors can directly identify what needs to be done and how to apply the
action.'' & Unsupported inference & The comment supplies no experimental design
or theorem statement. \\
``Therefore, this meets the criteria for a score of 5.'' & Score mapping error &
The preceding reasoning supports an explicit but vague action. \\
\bottomrule
\end{tabularx}
\captionsetup{hypcap=false}
\captionof{table}{Sentence-level error example. The unanimous consensus is
\texttt{supports\_wrong\_score}.}
\label{tab:rationale-sentence-example}
\end{minipage}
\vspace{0.8\baselineskip}
\begin{minipage}{\textwidth}
\normalsize\linespread{1.05}\selectfont
\setlength{\tabcolsep}{7pt}
\renewcommand{\arraystretch}{1.12}
\begin{tabularx}{\linewidth}{@{}X|X@{}}
\toprule
\multicolumn{2}{@{}p{\dimexpr\linewidth-2\tabcolsep-\arrayrulewidth\relax}@{}}{%
\textbf{Task: Actionability} \quad $\vert$ \quad
\textbf{Evaluation: Rationale Correction \& Rescoring} \hfill
\textbf{Gold label $=5$}} \\
\midrule
\multicolumn{2}{@{}p{\dimexpr\linewidth-2\tabcolsep-\arrayrulewidth\relax}@{}}{%
\textbf{DS prediction: 5 \,(\textcolor{green!50!black}{correct})}
\hfill
\textbf{NatRF prediction: 3 \,(\textcolor{red!70!black}{severe error})}} \\
\midrule
\multicolumn{2}{@{}>{\raggedright\arraybackslash}p{\dimexpr\linewidth-2\tabcolsep-\arrayrulewidth\relax}@{}}{%
\textbf{Evaluation input:} \textit{The training time of supernet should also be
listed for comparison.}} \\
\midrule
\rowcolor{gray!12}
\multicolumn{1}{@{}p{\dimexpr 0.5\linewidth-1.5\tabcolsep-\arrayrulewidth\relax}@{}}{%
\textbf{Original rationale}} &
\multicolumn{1}{@{}p{\dimexpr 0.5\linewidth-1.5\tabcolsep-\arrayrulewidth\relax}@{}}{%
\textbf{Corrected rationale}} \\
The request is a direct instruction, but it is
\textcolor{red!70!black}{judged explicit yet vague because it does not specify
where to list training time, which format to use, or how to compute it.} &
\textcolor{blue!70!black}{The requested addition is itself concrete and
executable: report the supernet's training time for comparison. The rubric
concerns whether the modification is specific, not presentation details such
as table placement or formatting.} \\
\midrule
\multicolumn{2}{@{}p{\dimexpr\linewidth-2\tabcolsep-\arrayrulewidth\relax}@{}}{%
\textbf{Record:} \texttt{actionability\_test\_0295}. \hfill
\textbf{Review:} GLM edit; MiniMax approved.} \\
\end{tabularx}
\nopagebreak
\vspace{-2pt}
\begin{tabularx}{\linewidth}{@{}l>{\centering\arraybackslash}X@{}}
\midrule
 & \textbf{Rescored prediction} \\
\textbf{A: natural rationale and score (NatRF)} & $3$ \\
\textbf{B: fixed original rationale and score (FixOrig)} & $3$ \\
\rowcolor{gray!12}
\textbf{C: fixed corrected rationale and score (FixCorr)} & \textbf{$5$} \\
\midrule
\multicolumn{2}{@{}p{\dimexpr\linewidth-2\tabcolsep\relax}@{}}{%
\hfill \textbf{Result: B reproduces A; C recovers the gold label.}} \\
\bottomrule
\end{tabularx}
\captionsetup{hypcap=false}
\captionof{table}{Rationale correction and fixed-prefix rescoring. FixOrig
reproduces NatRF, whereas FixCorr recovers the gold label.}
\label{tab:rationale-correction-example}
\end{minipage}

\end{table*}
